% ! TeX program = lualatex
% The geometric reading of the derivative: a curve, a point a on the x-axis,
% and the tangent line there.  Used by the midterm review slides.
%
% The curve is 3 - (x-3)^2/4 on [0, 6], chosen so the tangent at a = 1 has
% slope exactly 1 and touches at (1, 2) -- whole numbers, so the picture can
% be read off without the slope looking accidental.  Ticks are off: this
% frame talks about a and f'(a) rather than about particular values.
\documentclass[multi=tikzpicture,crop]{standalone}

\input{../typography.tex.preamble}
\input{../tikz.tex.preamble}

\pgfplotsset{tangent sketch/.style={
    basic curve,
    enlargelimits=false,
    xmin=-0.8, xmax=6.8,
    ymin=-0.8, ymax=4.4,
    % a grid for scale, but no numbers on it: this frame talks about a and
    % f(a), so the only labelled ticks are those two
    xtick distance={1}, ytick distance={1},
    xticklabels={}, yticklabels={},
    extra x ticks={1}, extra x tick label={\(a\)},
    extra x tick style={tick style={attn, thick}, tick label style={attn}},
    extra y ticks={2}, extra y tick label={\(f(a)\)},
    extra y tick style={tick style={attn, thick}, tick label style={attn}},
    width={220pt},
  }}

\begin{document}

\begin{tikzpicture}
  \begin{axis}[tangent sketch]
    \addplot[very thick, domain=0:6, samples=100] {3 - (x-3)^2/4};
    \node[anchor=north] at (axis cs:4.3,0.9) {\(y = f(x)\)};

    % slope 1 at a = 1, so the tangent is y = x + 1
    \addplot[main, very thick, domain=-0.5:3] {x + 1};
    \node[main, anchor=west] at (axis cs:3.2,3.7) {slope \(= \frac{dy}{dx}\)};

    \draw[attn, dashed] (axis cs:1,0) -- (axis cs:1,2) -- (axis cs:0,2);
    \draw[attn, fill=attn] (axis cs:1,2) circle (0.07);
  \end{axis}
\end{tikzpicture}
\end{document}
